\section[Fate of a redox-sensitive micropollutant]{Fate of a redox-sensitive micropollutant}

In the next part or the exercise the fate of the pharmaceutical compound 
phenazone will be simulated and the impact of the seasonally 
changing redox zonation will be investigated. 









To include {\color{red} \textbf{Phenazone}} take the following steps 

\begin{itemize}

\item Activate the species {\color{red} \textbf{Phenazone}} 
in the ${\color{blue} Solutions}$ Tab.
Enter an aqueous concentration 
of 1.0 $\times$ $10^{-6}$ mol/L for Phenazone for  
${\color{blue} Solutions}$ 1-12 but not for
${\color{blue} Solution}$ 0, i.e., 
not for the background water composition.

\item Define the reaction rate parameters for Phenazone under the 
${\color{blue} Rates}$ Tab. First, activate the Phenazone kinetics 
by ticking the first (C) box, but not the second (IM) box.
The latter would make Phenazone an immobile species.
Add a value of 5.0 $\times$ $10^{-6}$ 
for parm1 (reaction rate constant) and 
4.0 $\times$ $10^{-4}$ for parm2. Finally add \\
${\color{blue} Phenazone\ -1.0\ C11H12N2O\ 1.0}$ as ${\color{blue} Formula}$ \\
for phenazone to define the stoichiometry of the degradation reaction.

\end{itemize}

In the simulations of \cite{greskowiak_2006} it showed that 
the redox sensitivity (i.e., oxygen availability) dominated 
over the temperature effects on the phenazone degradation rates.
In their study this behaviour resulted in significantly higher 
removal rates during the colder month of the year, 
when the aerobic zone of the aquifer was significantly larger. 

A similar behaviour can also be seen if the breakthrough behaviour of 
phenazone is visualised at the extraction well.    